\cotefile{156-7}
\scanpage{58}
\mspage{57}

Notons qu'on \deleted{\unsure{donc|doit}} doit avoir
\[
\begin{aligned}
&\deleted{\operatorname{omb}(X)\cap\operatorname{omb}(Y)\ne\varnothing}\\
&\operatorname{omb}(X)=\operatorname{omb}(Y)\qquad (=|\varphi(X)|=|\varphi(Y)|)
\end{aligned}
\]
et qu'a priori il \unsure{n'est} pas exclu que $L=\varnothing$.

\uline{Exemple} $X$ espace topologique. On prend \deleted{$\unsure{M}$}
\[
M=\text{simplexes \uline{affines} dans }X
\]
avec la relation $\leq$ habituelle. \deleted{$\ll$ \unsure{la relation}} (mais en respectant les structures affines), $\ll$ à l'\unsure{avenant} (sans exiger que $|\sigma|\cap|\sigma'|$ soit vide ou un sous-simplexe commun), et $\unsure{\gtrless|\lessgtr|\gtreqless|\lesseqgtr}$ itou ($\unsure{\simeq|\cong}$ \unsure{presque}), d'où $\unsure{\mathscr{F}|\mathscr{E}}$, satisfaisant toutes les conditions\inserted{ de (ii)}. Mais considérons, sur une partie de $X$ homéomorphe à $[0,1]$, deux structures \unsure{affines} \deleted{\unsure{deux}} différentes, d'où $X,Y\in M$ avec \inserted{\unsure{isom.}} $L=\partial X=\partial Y$, $\varphi(X)=\varphi(Y)$, mais $X\ne Y$. Donc ici
\[
\varphi(X)\gtrless\varphi(Y)\quad\text{et}\quad\varphi(X)\ll\varphi(Y),
\]
sans avoir ni
\[
X\gtrless Y\quad\text{ni}\quad X\ll Y.
\]

% margin note written sideways; slanted upward along the left edge
\margin{Plus précisément\newline $(i)\Rightarrow(ii)+\varphi\text{ injectif}\Rightarrow(i\,a,c)^\circ$}

% A vertical marginal stroke brackets the corollary.
\uline{Cor} La condition (i) \deleted{\unsure{équivaut à}} \inserted{implique} (ii) + régularité de $\varphi$. \deleted{Il suffit \unsure{de} d'avoir l'injectivité de $\varphi|M$,} \inserted{Mais l'inverse n'est} pas vrai.

Bien sûr, i implique l'injectivité de $\varphi$, reste à prouver que $(ii)+\varphi|M\ \text{inj.}\Rightarrow b,c$ \unsure{(}ça a été prouvé déjà).

Prouvons c), plutôt, continuons. On a vu (utilisant l'injectivité de $\varphi$) que
